% ============================================================
% M3 成卷轮 S4 件2 —— 滚动测评卷（A）第二章（2.1–2.3）（16 题·40 分钟 100 分）
% 2026-09-14｜执行：M3 S4 卷件生产臂
% 骨架＝M2 滚动卷A main.tex 件型层 A~E 照抄；卷名/副题/键式改 M3（2章-滚A-N）。
% 卷式：单选 7×5＝35｜多选 2×5＝10（全对5/部分2/有错0）｜填空 3×5＝15｜解答 8+10+10+12＝40。
% 题源：蓝图 §二 16 席矩阵（扩容 13＋命制 2＋机动弹药 1）：
%   滚A-2 主钉 2章件6-#11／备钉 2章件12-#1 均「判断正误」三小条不合单选槽——
%   按蓝图 §四.3 装配期机动弹药启用 2章件7-#15（✅净三门空·0.94 简·圆域同节点·单选原生），
%   越蓝图钉位启用已登生产报告供复核（禁私造红线未触：源题全字照录）。
% 案B 附卷制：答案迁卷末附卷、% ans: 锚、速查表纯题档吞、\anskey 补偿发射。
% 形态转换（生产报告留痕）：滚A-6/滚A-14/滚A-15/滚A-16 设问微调（去选项/空线改「求/则…为」）；
%   滚A-8 命题列举开放题→五选项多选封装（命题文字逐字迁入选项）；滚A-10 结构不良保留原形；
%   滚A-11 单选→填空（去选项留空）。
% 编译门：xelatex main.tex ×2 / main-pure.tex ×2（sishou remember picture 硬依赖）。
% ============================================================
\documentclass[fontset=none]{ctexart}
\usepackage{qp-m3}
\ansblockgrayfalse   % 换装（测评/滚动＝8 开三栏件）：答案块一律括线模（方案草案§1.2.5/§5.1）
\providecommand{\ov}[1]{\overrightarrow{#1}}

% ------------------------------------------------------------
% 件型层 A｜页面：8 开横放三栏（承测评卷件型层）
% ------------------------------------------------------------
\geometry{paperwidth=399.8mm,paperheight=284.2mm,left=8.9mm,right=8.9mm,
  top=12.0mm,bottom=17.6mm,twoside,headheight=0pt,headsep=0pt,footskip=7.1mm}
\setlength{\topskip}{0pt}
\raggedbottom
\renewcommand{\normalsize}{\fontsize{10.5pt}{17.7pt}\selectfont}
\linespread{1}\normalsize
\setlength{\parindent}{0pt}
\setlength{\parskip}{0pt}
\newlength{\jpcolw}\setlength{\jpcolw}{120.0mm}
\newlength{\jpgap}\setlength{\jpgap}{11.0mm}
\xeCJKsetup{CJKglue={\hskip 0.04em plus 0.04em minus 0.01em}}

\definecolor{silklight}{gray}{0.8235}
\definecolor{silkdark}{gray}{0.6510}
\definecolor{juanrule}{gray}{0.1255}
\definecolor{subgrey}{gray}{0.4}

% ------------------------------------------------------------
% 件型层 B｜JT-01 丝带角饰（照抄测评卷件型层）
% ------------------------------------------------------------
\newcommand{\juanmathSize}{17pt}
\newcommand{\sishou}{%
  \begin{tikzpicture}[overlay,remember picture,x=1mm,y=1mm]
    \path (current page.north west) coordinate (O);
    \begin{scope}[shift={(O)}]
      \fill[silklight] (6.9,-9.39) -- (40.1,-9.39) -- (6.9,-42.59) -- cycle;
      \draw[white,line width=0.28mm] (9.21,-40.29) -- (9.21,-11.63) -- (37.87,-11.63);
      \fill[silkdark] (22.94,-8.21) -- (34.15,-8.21) -- (5.78,-36.58) -- (5.78,-25.37) -- cycle;
      \fill[silkdark] (34.15,-8.21) -- (34.15,-9.20) -- (33.16,-9.20) -- cycle;
      \fill[silkdark] (5.78,-36.58) -- (6.78,-36.58) -- (6.78,-35.58) -- cycle;
      \node[white,inner sep=0pt,anchor=center,rotate=45] at (17.23,-19.74)
        {\fontsize{\juanmathSize}{\juanmathSize}\selectfont\heitiao 数学};
    \end{scope}
  \end{tikzpicture}}

% ------------------------------------------------------------
% 件型层 C｜卷头零件＋题区零件（承测评卷件型层；卷名/副题/时间按滚A 卷式改文）
% ------------------------------------------------------------
\newcommand{\zeroheight}[1]{\raisebox{0pt}[0pt][0pt]{#1}}
\newcommand{\jpcen}[1]{\hbox to\jpcolw{\hfil\zeroheight{#1}\hfil}\nointerlineskip}
\newcommand{\jpright}[1]{\hbox to\jpcolw{\hfill\zeroheight{#1}}\nointerlineskip}
\newcommand{\vgt}[1]{\par\vskip\dimexpr#1-6.25mm\relax}

\newcommand{\tianxieg}[1]{{\fontsize{\qpJTbLabelSize}{13pt}\selectfont\heitiao #1：}%
  \rule[-0.62mm]{\qpJTbLineLen}{0.2mm}}
\newcommand{\jtianxie}{\tianxieg{班级}\hspace{\qpJTbGapGrp}\tianxieg{姓名}%
  \hspace{\qpJTbGapGrp}\tianxieg{得分}}
\newcommand{\juanmingSize}{20.5pt}
\newcommand{\jjuanming}{{\fontsize{\juanmingSize}{26pt}\selectfont\heizhang
  \renewcommand{\CJKglue}{\hskip 0pt}滚动测评卷（A）}}
\newcommand{\jjunrule}{{\color{juanrule}\rule{58.0mm}{\qpJTcRule}}}
\newcommand{\jfuti}{{\fontsize{\qpJTdSize}{18pt}\selectfont\heitiao\color{subgrey}%
  \renewcommand{\CJKglue}{\hskip 0pt}第二章（2.1–2.3）}}
\newcommand{\jptime}{{\fontsize{\qpJTeSize}{17.7pt}\selectfont（时间：40分钟\quad 分值：100分）}}
\newcommand{\zuhang}[2]{\par\noindent\hangindent=5.7mm\hangafter=1
  {\fontsize{\qpJTfSize}{17.7pt}\selectfont\heitiao #1}#2\par}
\newcommand{\ti}[2]{\par\noindent\hangindent=5.7mm\hangafter=1
  {\fontsize{11.4pt}{17.7pt}\selectfont\heihao{\numboldjian #1}.}\hspace{2.7mm}#2\par}
\newcommand{\fenzhi}[1]{（#1分）}
\newcommand{\optline}[1]{\par\noindent\hangindent=5.7mm\hangafter=0
  {\fontsize{10.5pt}{19pt}\selectfont #1}\par}
\newcommand{\optII}[2]{\makebox[53.05mm][l]{#1}#2}
\newcommand{\optIV}[4]{\makebox[21.5mm][l]{#1}\makebox[21.5mm][l]{#2}\makebox[21.5mm][l]{#3}#4}
\newcommand{\kw}{\hfill（\hspace{3.6mm}）\hspace*{5.4mm}}
\newcommand{\qpart}[1]{\par\noindent\hangindent=5.7mm\hangafter=0
  \hspace*{5.7mm}{\fontsize{10.5pt}{17.7pt}\selectfont #1}\par}

\newdimen\hA \hA=3.03mm
\newdimen\hB \hB=12.12mm
\newdimen\hC \hC=1.89mm
\newdimen\hD \hD=5.95mm
\newdimen\hE \hE=8.04mm
\newdimen\hF \hF=5.70mm
\newdimen\hG \hG=6.35mm
\newcommand{\jphead}{\hbox to\jpcolw{\sishou\hfill}\nointerlineskip
  \vskip\hA\jpright{\jtianxie}%
  \vskip\hB\jpcen{\jjuanming}%
  \vskip\hC\jpcen{\jjunrule}%
  \vskip\hD\jpcen{\jfuti}%
  \vskip\hE\jpcen{\jptime}%
  \vskip\hF
  \zuhang{一、选择题}{：本题共7小题，每小题5分，共35分．\ 在每小题给出的四个选项中，只有一项是符合题目要求的．}%
  \vgt{\hG}}

% ------------------------------------------------------------
% 件型层 D｜YJ-02 文字行页脚（同测评卷；卷名改滚A）
% ------------------------------------------------------------
\newcommand{\jpnum}[1]{{\fontsize{\qpYJbNumSize}{\qpYJbNumSize}\selectfont\numbold #1}}
\newcommand{\jptxtsize}{7.5pt}
\newcommand{\jptxt}{\fontsize{\jptxtsize}{\jptxtsize}\selectfont}
\newcommand{\juanname}{滚动测评卷（A）}
\newcommand{\pqt}{\ifnum\value{page}<10 0\fi\thepage}
\setlength{\headwidth}{\textwidth}
\fancyfoot[RO]{\jptxt\juanname\hspace{3.95mm}{\heitiao 滚动卷}\hspace{3.88mm}卷\hspace{2.25mm}\jpnum{\pqt}}
\fancyfoot[LE]{\jptxt\jpnum{\pqt}\hspace{1.9mm}卷\hspace{3.95mm}{\heitiao 羿郭工作室}%
  \hspace{3.0mm}高中数学\hspace{2.75mm}选择性必修第一册\hspace{2.8mm}RJB}

% ------------------------------------------------------------
% 件型层 E｜三栏排布器＋卷末答案速查表
% ------------------------------------------------------------
\newcommand{\jpcol}[1]{\raisebox{0pt}[0pt][0pt]{\vtop{\hsize\jpcolw\parskip0pt\parindent0pt
  \setlength\linewidth{\jpcolw}\setlength\columnwidth{\jpcolw}   % 换装补钉：qp-m3 括线盒按 \linewidth 取宽
  \fontsize{10.5pt}{17.7pt}\selectfont\relax#1}}}
\newcommand{\jpthree}[3]{\nointerlineskip\noindent\jpcol{#1}\hspace{\jpgap}\jpcol{#2}%
  \hspace{\jpgap}\jpcol{#3}\par}
\newcommand{\jplead}{\hbox{}\nointerlineskip}

\newcommand{\datu}[1]{{\fontsize{9pt}{12pt}\selectfont #1}}
\newcommand{\dabiao}{%
  \par\vgt{\hG}\noindent{\fontsize{10.5pt}{17.7pt}\selectfont\heitiao 答案速查}\par\nopagebreak
  \vskip 1.2mm\noindent\datu{\setlength{\tabcolsep}{1.4pt}%
  \begin{tabular}{|c|*{16}{c|}}
  \hline
  题号 & 1 & 2 & 3 & 4 & 5 & 6 & 7 & 8 & 9 & 10 & 11 & 12 & 13 & 14 & 15 & 16 \\ \hline
  答案 & A & B & B & C & B & A & D & ACE & ACD & \(\left(x\!+\!\frac{1}{2}\right)^{2}\!\!+\!\left(y\!+\!1\right)^{2}\!=\!\frac{9}{4}\) & \(\frac{13}{5}\) & \(\frac{\sqrt{2}}{2}\) & 略 & 略 & 略 & 略 \\ \hline
  \end{tabular}}\par}

\begin{document}
%装配轮S6三本跨本连续页码（G7方案B）setcounter 注入位
\setcounter{page}{313}
% ===================== 第 1 页：栏1 卷头＋Q1–Q2；栏2 Q3–Q5；栏3 Q6–Q7＋组行二＋Q8 =====================
\jpthree{\jphead
  \ti{1}{经过两点 \(C(3,1)\)、\(D(-2,0)\) 的直线 \(l_{1}\) 与经过点 \(M(1,-4)\) 且斜率为 \(\dfrac{1}{5}\) 的直线 \(l_{2}\) 的位置关系为\kw}
  \optline{\optIV{A．平行；}{B．垂直；}{C．重合；}{D．无法确定．}}
  \ti{2}{已知圆 \((x-1)^{2}+y^{2}=4\) 内一点 \(P(2,1)\)，则过 \(P\) 点的最短弦所在的直线方程是\kw}
  \optline{\optII{A．\(x-y-1=0\)；}{B．\(x+y-3=0\)；}}
  \optline{\optII{C．\(x+y+3=0\)；}{D．\(x=2\)．}}
}{\jplead
  \ti{3}{下列条件中，使得 \(l_{1}\perp l_{2}\) 的是\kw}
  \optline{①\(l_{1}\) 的斜率为 \(-\dfrac{2}{3}\)，\(l_{2}\) 经过点 \(A(1,1)\)，\(B\left(0,-\dfrac{1}{2}\right)\)；}
  \optline{②\(l_{1}\) 的倾斜角为 \(45^{\circ}\)，\(l_{2}\) 经过点 \(P(-2,-1)\)，\(Q(3,-5)\)；}
  \optline{③\(l_{1}\) 经过点 \(M(1,0)\)，\(N(4,-5)\)，\(l_{2}\) 经过点 \(R(-6,0)\)，\(S(-1,3)\)．}
  \optline{\optIV{A．①②；}{B．①③；}{C．②③；}{D．①②③．}}
  \ti{4}{在直线 \(2x-3y+5=0\) 上求点 \(P\)，使点 \(P\) 到 \(A(2,3)\) 的距离为 \(\sqrt{13}\)，则 \(P\) 点坐标是\kw}
  \optline{\optII{A．\((5,5)\)；}{B．\((-1,1)\)；}}
  \optline{\optII{C．\((5,5)\) 或 \((-1,1)\)；}{D．\((5,5)\) 或 \((1,-1)\)．}}
  \ti{5}{若不同的两点 \(P(a,b)\) 与 \(Q(b-1,a+1)\) 关于直线 \(l\) 对称，则直线 \(l\) 的倾斜角为\kw}
  \optline{\optIV{A．\(135^{\circ}\)；}{B．\(45^{\circ}\)；}{C．\(30^{\circ}\)；}{D．\(60^{\circ}\)．}}
}{\jplead
  \ti{6}{已知函数 \(y=x^{2}-2x-3\) 与 \(x\) 轴交于 \(A\)，\(B\) 两点，与 \(y\) 轴交于点 \(C\)，则 \(\triangle ABC\) 外接圆的方程为\kw}
  \optline{\optII{A．\(x^{2}+y^{2}-2x+2y-3=0\)；}{B．\(x^{2}+y^{2}-2x-2y-3=0\)；}}
  \optline{\optII{C．\(x^{2}+y^{2}+2x-2y-3=0\)；}{D．\(x^{2}+y^{2}-2x+2y+3=0\)．}}
  \ti{7}{若圆 \(C\colon x^{2}+y^{2}=9\) 上恰有 \(3\) 个点到直线 \(l\colon x-y+b=0\)（\(b>0\)）的距离为 \(2\)，\(l_{1}\colon x-y+3\sqrt{2}=0\)，则 \(l\) 与 \(l_{1}\) 间的距离为\kw}
  \optline{\optIV{A．\(1\)；}{B．\(\sqrt{2}\)；}{C．\(3\)；}{D．\(2\)．}}
  \vgt{\hG}
  \zuhang{二、选择题}{：本题共2小题，每小题5分，共10分．\ 在每小题给出的选项中，有多项符合题目要求．（全部选对得5分，部分选对得2分，有选错的得0分）}%
  \vgt{\hG}
  \ti{8}{在平面直角坐标系中，如果 \(x\) 与 \(y\) 都是整数，就称点 \((x,y)\) 为整点，则下列命题中正确的是\kw}
  \optline{A．存在这样的直线，既不与坐标轴平行又不经过任何整点；}
  \optline{B．如果 \(k\) 与 \(b\) 都是无理数，则直线 \(y=kx+b\) 不经过任何整点；}
  \optline{C．直线 \(l\) 经过无穷多个整点，当且仅当 \(l\) 经过两个不同的整点；}
  \optline{D．直线 \(y=kx+b\) 经过无穷多个整点的充分必要条件是：\(k\) 与 \(b\) 都是有理数；}
  \optline{E．存在恰经过一个整点的直线．}
}
\newpage
% ===================== 第 2 页：栏1 Q9＋组行三＋Q10–Q11；栏2 Q12＋组行四＋Q13–Q14；栏3 Q15–Q16＋速查表 =====================
\jpthree{\jplead
  \ti{9}{已知直线 \(l_{1}\colon 2x-y+1=0\) 与 \(l_{2}\colon 2x-y-9=0\)，点 \(A(1,2)\)、\(B(3,1)\)，则下列结论正确的是\kw}
  \optline{A．\(l_{1}\) 与 \(l_{2}\) 之间的距离为 \(2\sqrt{5}\)；}
  \optline{B．点 \(A\) 到直线 \(l_{1}\) 的距离为 \(\sqrt{5}\)；}
  \optline{C．点 \(B\) 到直线 \(l_{2}\) 的距离为 \(\dfrac{4\sqrt{5}}{5}\)；}
  \optline{D．过点 \(A\) 且与 \(l_{1}\) 垂直的直线方程为 \(x+2y-5=0\)．}
  \vgt{\hG}
  \zuhang{三、填空题}{：本题共3小题，每小题5分，共15分．}%
  \vgt{\hG}
  \ti{10}{在①经过直线 \(l_{1}\colon x-2y=0\) 与直线 \(l_{2}\colon 2x+y-1=0\) 的交点．②圆心在直线 \(2x-y=0\) 上．③被 \(y\) 轴截得弦长 \(|AB|=2\sqrt{2}\)；这三个条件中任选一个，补充下面问题中，若问题中的圆存在，求圆的方程；若问题中圆不存在，请说明理由．问题：是否存在圆 \(Q\)，点 \(A(-2,-1)\)，\(B(1,-1)\) 均在圆上，且圆 \(Q\)\_\_\_\_？}
  \ti{11}{直线 \(l_{1}\colon 3ax-y-2=0\) 和直线 \(l_{2}\colon (2a-1)x+5ay-1=0\) 分别过定点 \(A\) 和 \(B\)，则 \(|AB|=\)\kongbai}
}{\jplead
  \ti{12}{过点 \(P(1,\sqrt{2})\) 的直线 \(l\) 将圆 \(C\colon (x-2)^{2}+y^{2}=4\) 分成两段弧，当劣弧所对的圆心角最小时，直线 \(l\) 的斜率 \(k\) 为\kongbai}
  \vgt{\hG}
  \zuhang{四、解答题}{：本题共4小题，共40分．\ 解答应写出文字说明、证明过程或演算步骤．}%
  \vgt{\hG}
  \ti{13}{\fenzhi{8}已知直线 \(l\) 经过点 \(A(1,-2)\)，\(B(-3,4)\)．}
  \qpart{（1）求直线 \(l\) 的方程；}
  \qpart{（2）判断 \(l\) 与直线 \(m\colon 6x+4y-7=0\) 的位置关系，并说明理由．}
  \ti{14}{\fenzhi{10}在平面直角坐标系 \(xOy\) 中，\(A\) 为直线 \(l\colon y=2x\) 上在第一象限内的点，\(B(5,0)\)，以 \(AB\) 为直径的圆 \(C\) 与直线 \(l\) 交于另一点 \(D\)．若 \(\overrightarrow{AB}\cdot\overrightarrow{CD}=0\)，求点 \(A\) 的横坐标的值．}
}{\jplead
  \ti{15}{\fenzhi{10}已知点 \(A(-1,0)\)，\(B(1,0)\)，若圆 \((x-2a+1)^{2}+(y-2a-2)^{2}=1\) 上存在点 \(M\) 满足 \(\overrightarrow{MA}\cdot\overrightarrow{MB}=3\)，求满足条件的实数 \(a\) 的取值范围．}
  \ti{16}{\fenzhi{12}一座圆拱桥，当水面在某位置时，拱顶离水面 \(2\,\mathrm{m}\)，水面宽 \(12\,\mathrm{m}\)，当水面下降 \(1\,\mathrm{m}\) 后，求此时的水面宽．}
  \ifshowans\dabiao\fi   % 换装案B：速查表＝答案承载件，纯题档随开关吞（防漏答案）
  \ifshowans\else % 案B双锚补丁：附卷页整页跳过，纯题档在此补偿发射 16 键（两档 M3-ANSKEY 恒等，对号门两档可跑）
  \anskey{2章-滚A-1}\anskey{2章-滚A-2}\anskey{2章-滚A-3}\anskey{2章-滚A-4}\anskey{2章-滚A-5}\anskey{2章-滚A-6}\anskey{2章-滚A-7}\anskey{2章-滚A-8}\anskey{2章-滚A-9}\anskey{2章-滚A-10}\anskey{2章-滚A-11}\anskey{2章-滚A-12}\anskey{2章-滚A-13}\anskey{2章-滚A-14}\anskey{2章-滚A-15}\anskey{2章-滚A-16}
  \fi
}
% ===================== 卷末附卷：参考答案（案B 附卷制：仅含详解档成页；纯题档整页跳过） =====================
\ifshowans
\newpage
\jpthree{\jplead
  \par\vskip 2.0mm\noindent{\fontsize{\qpJTfSize}{17.7pt}\selectfont\heitiao 参考答案}\hspace{0.6em}{\fontsize{\qpJTeSize}{17.7pt}\selectfont（滚动测评卷（A）·含详解档）}\par\nopagebreak
  \begin{ansblock}[2章-滚A-1]
  % ans:2章-滚A-1
  \ansitem{1}{A}
  \end{ansblock}
  \begin{ansblock}[2章-滚A-2]
  % ans:2章-滚A-2
  \ansitem{2}{B}
  \end{ansblock}
  \begin{ansblock}[2章-滚A-3]
  % ans:2章-滚A-3
  \ansitem{3}{B（①③）}
  \end{ansblock}
  \begin{ansblock}[2章-滚A-4]
  % ans:2章-滚A-4
  \ansitem{4}{C}
  \end{ansblock}
  \begin{ansblock}[2章-滚A-5]
  % ans:2章-滚A-5
  \ansitem{5}{B（\(45^{\circ}\)）}
  \end{ansblock}
  \begin{ansblock}[2章-滚A-6]
  % ans:2章-滚A-6
  \ansitem{6}{A（\(x^{2}+y^{2}-2x+2y-3=0\)）}
  \end{ansblock}
  \begin{ansblock}[2章-滚A-7]
  % ans:2章-滚A-7
  \ansitem{7}{D（\(2\)）}
  \end{ansblock}
}{\jplead
  \begin{ansblock}[2章-滚A-8]
  % ans:2章-滚A-8
  \ansitem{8}{ACE}
  \end{ansblock}
  \begin{ansblock}[2章-滚A-9]
  % ans:2章-滚A-9
  \ansitem{9}{ACD}
  \ansnote{解析}{\(l_{1}\parallel l_{2}\)。A：\(d=\dfrac{|1-(-9)|}{\sqrt{2^{2}+(-1)^{2}}}=\dfrac{10}{\sqrt{5}}=2\sqrt{5}\)，正确；B：\(A\) 代入 \(l_{1}\) 左侧 \(=1\)，距离 \(\dfrac{\sqrt{5}}{5}\neq\sqrt{5}\)，错误；C：\(B\) 代入 \(l_{2}\) 左侧 \(=-4\)，距离 \(\dfrac{4\sqrt{5}}{5}\)，正确；D：\(l_{1}\) 斜率 \(2\)，垂线斜率 \(-\dfrac{1}{2}\)，\(y-2=-\dfrac{1}{2}(x-1)\) 即 \(x+2y-5=0\)，正确。故选 ACD。}
  \end{ansblock}
  \begin{ansblock}[2章-滚A-10]
  % ans:2章-滚A-10
  \ansitem{10}{\(\left(x+\dfrac{1}{2}\right)^{2}+(y+1)^{2}=\dfrac{9}{4}\)（三条件任选，圆均存在）}
  \end{ansblock}
  \begin{ansblock}[2章-滚A-11]
  % ans:2章-滚A-11
  \ansitem{11}{\(\dfrac{13}{5}\)}
  \end{ansblock}
  \begin{ansblock}[2章-滚A-12]
  % ans:2章-滚A-12
  \ansitem{12}{\(\dfrac{\sqrt{2}}{2}\)}
  \end{ansblock}
}{\jplead
  \begin{ansblock}[2章-滚A-13]
  % ans:2章-滚A-13
  \ansitem{13}{(1)\(3x+2y+1=0\)；(2)\(l\parallel m\)（平行）}
  \ansnote{解析}{(1) \(k=\dfrac{4-(-2)}{-3-1}=-\dfrac{3}{2}\)，\(y-4=-\dfrac{3}{2}(x+3)\)，整理得 \(3x+2y+1=0\)，经验证 \(A\)、\(B\) 均满足。(2) 一次项系数比 \(\dfrac{3}{6}=\dfrac{2}{4}=\dfrac{1}{2}\)，而常数比 \(\dfrac{1}{-7}\neq\dfrac{1}{2}\)，故 \(l\parallel m\)。}
  \end{ansblock}
  \begin{ansblock}[2章-滚A-14]
  % ans:2章-滚A-14
  \ansitem{14}{\(3\)}
  \ansnote{解析}{设 \(A(a,2a)\)（\(a>0\)），圆心 \(C\left(\dfrac{a+5}{2},a\right)\)，\(\odot C\)：\((x-5)(x-a)+y(y-2a)=0\)，与 \(y=2x\) 联立解得 \(D(1,2)\)。\(\overrightarrow{AB}=(5-a,-2a)\)，\(\overrightarrow{CD}=\left(1-\dfrac{a+5}{2},2-a\right)\)，由 \(\overrightarrow{AB}\cdot\overrightarrow{CD}=0\) 得 \(a^{2}-2a-3=0\)，\(a=3\)（\(a=-1\) 舍），故点 \(A\) 的横坐标为 \(3\)。}
  \end{ansblock}
  \begin{ansblock}[2章-滚A-15]
  % ans:2章-滚A-15
  \ansitem{15}{\(-1\leq a\leq\dfrac{1}{2}\)}
  \ansnote{解析}{设 \(M(x,y)\)：\(\overrightarrow{MA}\cdot\overrightarrow{MB}=(-x-1)(-x+1)+y^{2}=3\)，得 \(M\) 的轨迹 \(x^{2}+y^{2}=4\)。圆 \((x-2a+1)^{2}+(y-2a-2)^{2}=1\) 圆心 \((2a-1,2a+2)\)、半径 \(1\)，与 \(x^{2}+y^{2}=4\) 有公共点 \(\iff 1\leq\sqrt{(2a-1)^{2}+(2a+2)^{2}}\leq 3\)，即 \(8a^{2}+4a-4\leq 0\)，解得 \(-1\leq a\leq\dfrac{1}{2}\)。故满足条件的实数 \(a\) 的取值范围为 \(\left[-1,\dfrac{1}{2}\right]\)。}
  \end{ansblock}
  \begin{ansblock}[2章-滚A-16]
  % ans:2章-滚A-16
  \ansitem{16}{\(2\sqrt{51}\,\mathrm{m}\)}
  \ansnote{解析}{以拱顶为原点、过拱顶的水平切线为 \(x\) 轴建系，设圆 \(x^{2}+(y+r)^{2}=r^{2}\)。由拱顶离水面 \(2\,\mathrm{m}\)、水面宽 \(12\,\mathrm{m}\)，点 \((6,-2)\) 在圆上：\(36+64=100=r^{2}\)，\(r=10\)，圆 \(x^{2}+(y+10)^{2}=100\)。水面下降 \(1\,\mathrm{m}\) 后 \(y=-3\)：\(x^{2}=100-49=51\)，半宽 \(\sqrt{51}\)，水面宽 \(2\sqrt{51}\,\mathrm{m}\)。}
  \end{ansblock}
}
\fi
\end{document}
